\documentclass[11pt]{article}
\usepackage[localtheorems]{raeez-math-template}

\newtheorem{theorem}{Theorem}
\newtheorem{proposition}[theorem]{Proposition}
\newtheorem{lemma}[theorem]{Lemma}
\newtheorem{conjecture}[theorem]{Conjecture}
\newtheorem{definition}[theorem]{Definition}
\newtheorem{remark}[theorem]{Remark}

\providecommand{\fg}{\mathfrak{g}}
\providecommand{\fh}{\mathfrak{h}}
\providecommand{\CoHA}{\mathrm{CoHA}}
\providecommand{\Coh}{\mathrm{Coh}}
\providecommand{\Rep}{\mathrm{Rep}}
\providecommand{\End}{\mathrm{End}}
\providecommand{\Hilb}{\mathrm{Hilb}}
\providecommand{\Stab}{\mathrm{Stab}}
\providecommand{\MO}{\mathrm{MO}}
\providecommand{\Perf}{\mathrm{Perf}}
\providecommand{\CC}{\mathbb{C}}
\providecommand{\ZZ}{\mathbb{Z}}
\providecommand{\PP}{\mathbb{P}}
\providecommand{\cO}{\mathcal{O}}
\providecommand{\cZ}{\mathcal{Z}}

\title{CY-C Three Routes for $K3 \times E$:\\
CoHA, Koszul Duality, and Stable Envelopes}
\author{Raeez Lorgat}
\date{2026-04-22}

\begin{document}
\maketitle

\begin{abstract}
Let $S$ be a smooth projective K3 surface, let $E$ be an elliptic
curve, and set $X=S\times E$. The CY-C comparison asks whether the
compact Hall positive half, the Koszul-dual chiral category, and the
Maulik-Okounkov stable-envelope algebra assemble to one completed
Hall-Drinfeld and Borcherds target. The comparison is conditional. The
compact positive half, non-degenerate Hopf pairing, completion, Koszul
model, stable-envelope fixed-locus geometry, and comparison maps are
part of the construction data. With those data supplied, the three
routes define a comparison problem on the Mukai and equivariant loci at
generic $q$. Without those data, CY-C remains conjectural and the
global quantum vertex chiral group $G(X)$ is unconstructed in general.
For $K3\times E$ the numerical invariants are fixed:
$\kappa_{\mathrm{cat}}(X)=0$, $\kappa_{\mathrm{ch}}(X)=0$ on the
compact total-space Hodge reading,
$\kappa_{\mathrm{ch}}^{\mathrm{Heis}}(X)=3$,
$\kappa_{\mathrm{BKM}}(\Delta_5)=5$, and
$\kappa_{\mathrm{fiber}}(X)=24$.
\end{abstract}

\tableofcontents

\section{Objects and Status}
\label{sec:objects-status}

The compact Calabi-Yau threefold is always
\[
 X=S\times E.
\]
An auxiliary curve $\Sigma_g$ is a chiral test curve or a stage-$2$
specialisation curve. Only at $g=1$ does this curve coincide with the
elliptic factor $E$ in the Calabi-Yau threefold.

\begin{definition}[The conditional target $G(X)$]
\label{def:conditional-GX}
A quantum vertex chiral group
\[
 G(X):=D_\hbar^\wedge(Y^+(X))
\]
is defined only after the following package is supplied:
\begin{enumerate}[label=\textup{(H\arabic*)}, leftmargin=2.2em]
\item a positive half $Y^+(X)$ in a completed $E_1$-chiral category;
\item a completed negative half $Y^-(X)$;
\item a non-degenerate Hopf pairing
$Y^+(X)\widehat\otimes Y^-(X)\to\mathbf{1}$;
\item a completed Drinfeld double $D_\hbar^\wedge(Y^+(X))$;
\item a centre comparison
\[
 \cZ\bigl(\Rep^{E_1}(Y^+(X))\bigr)
 \simeq
 \Rep^{E_2}\bigl(D_\hbar^\wedge(Y^+(X))\bigr);
\]
\item route comparison maps preserving the framed stage-$2$
specialisation, root grading, coproduct, pairing, and completion.
\end{enumerate}
For a general compact CY$_3$ this package is absent. Thus
$G(X)$ is unconstructed in general.
\end{definition}

\begin{definition}[Three route outputs]
\label{def:three-route-outputs}
The three constructions considered here are:
\begin{enumerate}[label=\textup{(\roman*)}, leftmargin=2.0em]
\item the compact Hall route
\[
 C^{\mathrm{Hall}}(X)
 :=
 D_\hbar^\wedge\bigl(Y^+_{\mathrm{Hall}}(X)\bigr),
\]
defined after compact critical CoHA correspondences, a positive half,
a non-degenerate Hopf pairing, and completion are supplied;
\item the Koszul route
\[
 C^{\mathrm{Kos}}(X,\Sigma_g)
 :=
 \End^\otimes\!\bigl(\Rep(A_{X,\Sigma_g}^{!})\bigr),
\]
defined after the stage-$2$ chiral algebra $A_{X,\Sigma_g}$ is
constructed and a Koszul dual model $A_{X,\Sigma_g}^{!}$ is fixed;
\item the stable-envelope route
\[
 C^{\mathrm{MO}}(S,T)
 :=
 \langle R^{\mathrm{MO}}_{\mathfrak C}(z)\rangle,
\]
the algebra generated by Maulik-Okounkov stable-envelope $R$-matrices
on a symplectic resolution with a torus action $T$ and controlled fixed
loci.
\end{enumerate}
\end{definition}

\begin{remark}[K3 BKM branch and K3 Yangian branch]
\label{rem:bkm-yangian-separation}
The BKM-side object attached to $K3\times E$ is the Hall-Drinfeld and
Borcherds branch governed by $\Delta_5$. The Mukai self-mirror Yangian
branch is a separate K3 lattice construction. The abelian Heisenberg
locus witnesses both branches at their Cartan or Fock-space shadows.
Beyond those shadows, the full compact Hall double and the K3 Yangian
remain distinct.
\end{remark}

\section{Numerical Invariants}
\label{sec:numerical-invariants}

\begin{theorem}[Four subscripted values on $K3\times E$]
\label{thm:k3e-four-values}
Let $Y=K3\times E$. The values attached to the four distinct
constructions are
\[
\begin{array}{ccl}
\kappa_{\mathrm{cat}}(Y) &=& 0,\\[2pt]
\kappa_{\mathrm{ch}}(Y) &=& 0
  \quad\text{on the compact total-space Hodge reading},\\[2pt]
\kappa_{\mathrm{ch}}^{\mathrm{Heis}}(Y) &=& 3,\\[2pt]
\kappa_{\mathrm{BKM}}(\Delta_5) &=& 5,\\[2pt]
\kappa_{\mathrm{fiber}}(Y) &=& 24.
\end{array}
\]
The set $\{0,3,5,24\}$ records four construction sources rather than
four evaluations of one functor.
\end{theorem}

\begin{proof}
Kunneth gives
\[
 h^{0,\bullet}(K3\times E)=(1,1,1,1),
\]
because $h^{0,\bullet}(K3)=(1,0,1)$ and $h^{0,\bullet}(E)=(1,1)$.
Hence
\[
 \kappa_{\mathrm{cat}}(Y)
 =
 \chi(\cO_Y)
 =
 \chi(\cO_{K3})\chi(\cO_E)
 =
 2\cdot 0
 =
 0,
\]
and the compact Hodge supertrace is
\[
 \kappa_{\mathrm{ch}}(Y)
 =
 \sum_{q=0}^{3}(-1)^q h^{0,q}(Y)
 =
 1-1+1-1
 =
 0.
\]
The Heisenberg-Mukai specialisation is a different stage-$2$ reading:
it reduces to the rank-$3$ Lorentzian Heisenberg sector and gives
$\kappa_{\mathrm{ch}}^{\mathrm{Heis}}(Y)=3$. Borcherds' weight theorem
gives the weight of the multiplicative lift as one half of the
discriminant-zero coefficient. In Eichler-Zagier normalisation for the
primitive K3 input, $c_1(0)=10$, whence
\[
 \kappa_{\mathrm{BKM}}(\Delta_5)=c_1(0)/2=5.
\]
Finally
\[
 \kappa_{\mathrm{fiber}}(Y)
 =
 \operatorname{rk}\widetilde H(K3,\ZZ)
 =
 1+22+1
 =
 24.
\]
These four computations use different inputs: structure-sheaf Euler
characteristic, compact Hodge supertrace, Heisenberg specialisation,
Borcherds lift, and Mukai-lattice rank.
\end{proof}

\begin{proposition}[Primitive CHL constants]
\label{prop:primitive-chl-constants}
For the primitive CHL denominator ladder
$N\in\{1,2,3,4,6\}$,
\[
 \bigl(c_N(0)\bigr)=(10,8,6,4,2)
\]
and
\[
 \bigl(\kappa_{\mathrm{BKM}}(\Phi_N)\bigr)
 =
 \bigl(c_N(0)/2\bigr)
 =
 (5,4,3,2,1).
\]
The Gritsenko-Clery diagonal-divisor catalogue is a separate eight-row
catalogue with its own cover groups and weights.
\end{proposition}

\begin{proof}
Borcherds' weight theorem gives
$\operatorname{wt}\mathrm{BL}(\phi)=c(0)/2$ for the multiplicative
lift of a weak Jacobi or vector-valued input with discriminant-zero
coefficient $c(0)$. The primitive CHL inputs have the displayed
constant terms. Taking one half gives the displayed
$\kappa_{\mathrm{BKM}}(\Phi_N)$ tuple. The Gritsenko-Clery catalogue is
indexed by different diagonal-divisor data; importing the CHL tuple
there requires a row-by-row overlap theorem.
\end{proof}

\section{Conditional Three-Route Comparison}
\label{sec:conditional-comparison}

\begin{conjecture}[CY-C three-route comparison]
\label{conj:cy-c-three-route-comparison}
Fix $X=S\times E$ and generic $q\in\CC^\times$. Suppose that:
\begin{enumerate}[label=\textup{(C\arabic*)}, leftmargin=2.2em]
\item the compact Hall positive half $Y^+_{\mathrm{Hall}}(X)$ is
constructed with coproduct and completion;
\item the Serre-dual negative half and non-degenerate Hopf pairing are
constructed;
\item the stage-$2$ chiral algebra $A_{X,\Sigma_g}$ has a Koszul model
$A_{X,\Sigma_g}^{!}$;
\item the Maulik-Okounkov construction is available on a torus
equivariant symplectic-resolution chart attached to the same Mukai
data;
\item the comparison maps preserve the root grading, PBW filtration,
Serre kernel, coproduct, Hopf pairing, centre, and completion.
\end{enumerate}
Then the three route outputs are isomorphic in the completed
Hall-Drinfeld target:
\[
 C^{\mathrm{Hall}}(X)
 \simeq
 C^{\mathrm{Kos}}(X,\Sigma_g)
 \simeq
 C^{\mathrm{MO}}(S,T)
 \simeq
 D_\hbar^\wedge(Y^+_{\mathrm{Hall}}(X)).
\]
\end{conjecture}

\begin{remark}[Conjectural status]
\label{rem:why-conjecture}
The scalar identities above supply numerical constraints only. The
route maps in Conjecture~\ref{conj:cy-c-three-route-comparison}, the
compact Hall pairing, the Koszul dual model, and the stable-envelope
fixed-locus input remain algebraic construction data. CY-C remains
conjectural because those data are independent of the Borcherds
denominator and Kunneth.
\end{remark}

\begin{proposition}[Six routes are distinct constructions]
\label{prop:six-routes-distinct}
The three routes in Definition~\ref{def:three-route-outputs} sit inside
the six-route comparison for the conjectural target $G(K3\times E)$.
Only the framed CY-to-chiral route applies
\[
 \Phi_3=\mathrm{Sp}^{\mathrm{ch}}_{\Sigma_2,C}\circ\Phi_3^{\mathrm{FA}}
\]
to $\Perf(K3\times E)$. The Borcherds, lattice-VOA, Kummer,
sigma-model, and Schur-sector routes are independent constructions.
Their convergence is a CY-C assertion, not a consequence of repeating
$\Phi_3$ six times.
\end{proposition}

\section{Defined Loci}
\label{sec:defined-loci}

\begin{proposition}[Generic $q$, Heisenberg locus]
\label{prop:g0-heisenberg-locus}
On the Heisenberg or Kummer locus, Nakajima-Grojnowski theory supplies
the Fock representation
\[
 \bigoplus_{n\geq 0} H^*(\Hilb^n S)
 \cong
 \mathcal{F}_{\widetilde H(S,\ZZ)}.
\]
There the Hall and stable-envelope constructions have a common
Heisenberg shadow. The full equality with the Koszul route requires a
constructed Koszul dual $A_{X,\Sigma_0}^{!}$ and comparison maps to the
same completed double.
\end{proposition}

\begin{proposition}[Elliptic factor and stable envelopes]
\label{prop:g1-stable-envelope-scope}
At $g=1$ the test curve is the elliptic factor $E$ and the geometry is
$X=S\times E$. The compact Hall route is conditional on compact
critical CoHA correspondences, purity sufficient for PBW, a
non-degenerate Hopf pairing, and completion. The stable-envelope route
requires a torus action on a symplectic-resolution chart with controlled
fixed loci. A generic K3 surface has no algebraic torus action of this
kind. A finite Mukai symplectic automorphism group supplies finite
equivariance data; Maulik-Okounkov stable envelopes require the
additional torus action.
\end{proposition}

\begin{proposition}[Schottky boundary]
\label{prop:schottky-boundary}
At genus $2$, $\dim \mathcal M_2=\dim \mathcal A_2=3$, and the Torelli
map identifies $\mathcal M_2$ with the open Jacobian locus in
$\mathcal A_2$, with the decomposable Humbert divisor as boundary. A
comparison between curve-side chiral data and Siegel automorphic data
therefore requires pullback along the Torelli map. For $g=3$ the
dimensions are again equal. For $g\geq 4$ the Schottky locus has
codimension
\[
 \frac{g(g+1)}{2}-(3g-3)>0
\]
inside $\mathcal A_g$. A genus-$g$ route comparison with $g\geq 3$
must specify Schottky-compatible automorphic input before it can be
stated.
\end{proposition}

\begin{proposition}[Roots of unity]
\label{prop:roots-of-unity-scope}
Let $q=\zeta_N=e^{2\pi i/N}$ with $N\in\{2,3,4,6\}$. The generic
comparison problem changes to a cyclotomic problem. The Hall route
requires a Lusztig or cyclotomic quotient of the positive half; the
Koszul route requires a fusion or semisimplified tensor category; the
stable-envelope route requires specialised $R$-matrix data. At $N=2$
the conditional K3 quantum-toroidal model has
\[
 p_{24}(2)=24+24+\binom{24}{2}=324
\]
simple objects, and its explicit $S$-matrix has full rank $324$ in the
model computation. This count is a root-of-unity datum. The generic
completed Hall double requires the cyclotomic quotient and comparison
functor before comparison with this datum.
\end{proposition}

\section{Tensor Tests}
\label{sec:tensor-tests}

\begin{proposition}[Boundary of Heisenberg tensor data]
\label{prop:heisenberg-tensor-boundary}
The Nakajima Heisenberg construction on the K3 fibre and the elliptic
oscillator data determine a Fock-space shadow. The compact Hall double
requires additional data: the Borcherds denominator character, Weyl
vector, imaginary-root multiplicities, Serre kernel, Hall bracket, Hopf
pairing, and completion. Consequently a tensor expression built only
from the K3 Heisenberg algebra and the elliptic oscillator algebra is
insufficient for CY-C.
\end{proposition}

\begin{proof}
Nakajima's construction uses the Mukai lattice and incidence
correspondences on Hilbert schemes of points on the surface. It gives a
Heisenberg action and its Fock representation. The BKM branch uses the
Borcherds product attached to $\Delta_5$, whose denominator identity
records root multiplicities and Weyl chamber data. The tensor product
of surface Heisenberg operators and elliptic oscillators lacks the
Borcherds product, Hall pairing, and Serre kernel. Hence the compact
Hall-Drinfeld double requires the additional comparison data listed in
Definition~\ref{def:conditional-GX}.
\end{proof}

\section{Separated Loci}
\label{sec:separated-loci}

\begin{center}
\renewcommand{\arraystretch}{1.25}
\begin{tabular}{llll}
\toprule
Locus & Hall route & Koszul route & Stable-envelope route \\
\midrule
Heisenberg or Kummer, generic $q$ &
Fock shadow &
requires $A^!$ &
equivariant chart \\
Generic K3, $g=1$, generic $q$ &
conditional compact double &
requires Koszul model &
no generic torus chart \\
Mukai equivariant locus, generic $q$ &
equivariant Hall data &
equivariant $A^!$ &
requires torus chart \\
Genus $2$ &
curve-side data &
curve-side $A^!$ &
pullback along Torelli \\
$g\geq 3$ &
Schottky-compatible input required &
Schottky-compatible input required &
Schottky-compatible input required \\
$q=\zeta_N$ &
cyclotomic quotient &
fusion category &
specialised $R$-matrix \\
\bottomrule
\end{tabular}
\end{center}

\section{References}
\label{sec:references}

\begin{thebibliography}{99}

\bibitem{Borcherds1998}
R. Borcherds,
\emph{Automorphic forms with singularities on Grassmannians},
Invent. Math. 132 (1998), 491-562.

\bibitem{GritsenkoNikulin1998}
V. Gritsenko and V. Nikulin,
\emph{Automorphic forms and Lorentzian Kac-Moody algebras. Part II},
Internat. J. Math. 9 (1998), 201-275.

\bibitem{Gritsenko1999}
V. Gritsenko,
\emph{Irrationality of the moduli spaces of polarized abelian surfaces},
Abh. Math. Sem. Univ. Hamburg 69 (1999), 93-117.

\bibitem{CleryGritsenko2011}
F. Clery and V. Gritsenko,
\emph{The Siegel modular forms of genus 2 with the simplest divisor},
Proc. Lond. Math. Soc. 102 (2011), 1024-1052.

\bibitem{Mukai1987}
S. Mukai,
\emph{On the moduli space of bundles on K3 surfaces. I},
Vector bundles on algebraic varieties, Tata Institute, 1987.

\bibitem{Huybrechts2016}
D. Huybrechts,
\emph{Lectures on K3 Surfaces},
Cambridge University Press, 2016.

\bibitem{Nakajima1997}
H. Nakajima,
\emph{Heisenberg algebra and Hilbert schemes of points on projective
surfaces},
Ann. of Math. 145 (1997), 379-388.

\bibitem{Grojnowski1996}
I. Grojnowski,
\emph{Instantons and affine algebras I: The Hilbert scheme and vertex
operators},
Math. Res. Lett. 3 (1996), 275-291.

\bibitem{SchiffmannVasserot2012}
O. Schiffmann and E. Vasserot,
\emph{The elliptic Hall algebra and the K-theory of the Hilbert scheme
of $\mathbb A^2$},
Duke Math. J. 162 (2013), 279-366.

\bibitem{MaulikOkounkov2012}
D. Maulik and A. Okounkov,
\emph{Quantum Groups and Quantum Cohomology},
Asterisque 408 (2019).

\bibitem{DavisonMeinhardt2020}
B. Davison and S. Meinhardt,
\emph{Cohomological Donaldson-Thomas theory of a quiver with potential
and quantum enveloping algebras},
Invent. Math. 221 (2020), 777-871.

\bibitem{OberdieckPandharipande2018}
G. Oberdieck and R. Pandharipande,
\emph{Curve counting on K3 times elliptic curves and the Igusa cusp
form},
J. Topol. 11 (2018), 915-943.

\bibitem{DMVV1997}
R. Dijkgraaf, G. Moore, E. Verlinde, and H. Verlinde,
\emph{Elliptic genera of symmetric products and second quantized
strings},
Comm. Math. Phys. 185 (1997), 197-209.

\end{thebibliography}

\end{document}
